\documentclass[a4paper,11pt]{article}
\usepackage[utf8]{inputenc}
\usepackage[T1]{fontenc}
\usepackage[french]{babel}
\usepackage{lmodern}
\usepackage{amsmath,amssymb,amsthm}
\usepackage{mathrsfs}
\usepackage{enumitem}
\usepackage{multicol}

\setlist[enumerate]{label=\textbf{\textcolor{Couleur8}{\arabic*.}}, left=1em}

% Si vous voulez aussi modifier les listes enumerate imbriquées :
\setlist[enumerate,2]{label=\textcolor{Couleur8}{\alph*)}}
\setlist[enumerate,3]{label=\textcolor{Couleur8}{\roman*)}}
\setlist[enumerate,4]{label=\textcolor{Couleur8}{\Alph*)}}
\usepackage{graphicx}
\usepackage{xcolor}
\usepackage{tcolorbox}
\usepackage{geometry}
\usepackage{fancyhdr}
\usepackage{titlesec}
\usepackage{cancel}
\usepackage{ifthen}
\usepackage{tikz}
\usetikzlibrary{arrows.meta}
\usepackage{tkz-tab}
\usepackage{eurosym}

% OPTION PROF/ÉLÈVE
% Décommentez UNE des deux lignes suivantes :
\newboolean{prof}\setboolean{prof}{true}   % Version professeur (avec corrections des devoirs)
%\newboolean{prof}\setboolean{prof}{false}  % Version élève (sans corrections des devoirs)

\definecolor{Couleur1}{HTML}{4A5D7A} %Th
\definecolor{Couleur2}{HTML}{A5B5D6} %Prop
\definecolor{Couleur3}{HTML}{A8C68A} %Méthode
\definecolor{Couleur6}{HTML}{8A9CC1} %Def
\definecolor{Couleur7}{HTML}{E6B459} %Rq Voc
\definecolor{Couleur8}{HTML}{D36740} %Titres
\definecolor{correction}{HTML}{D36740} % Couleur pour les corrections
\definecolor{coopmathscorr}{HTML}{F15929} % Couleur corrections coopmaths

% Configuration des marges
\geometry{
	top=2cm,
	bottom=2cm,
	inner=1.5cm,
	outer=1.5cm,
}

% Police
\renewcommand{\familydefault}{\sfdefault}

% Compteurs
\newcounter{seancenum}
\newcounter{seancenumD}
\setcounter{seancenum}{1}
\setcounter{seancenumD}{1}

% Configuration des en-têtes et pieds de page
\pagestyle{fancy}
\fancyhf{}
\ifthenelse{\boolean{prof}}{
    \fancyhead[L]{\textcolor{Couleur8}{\textbf{Corrections Automatismes - Classe de seconde - Version Professeur}}}
}{
    \fancyhead[L]{\textcolor{Couleur8}{\textbf{Corrections Automatismes - Séance 25 - Classe de seconde}}}
}
\fancyhead[R]{\textcolor{Couleur8}{\textbf{2025-2026}}}
\fancyfoot[L]{\textcolor{Couleur8}{Mme Bertail et Mme Pinto}}
\fancyfoot[R]{\textcolor{Couleur8}{\thepage}}
\renewcommand{\headrulewidth}{1pt}
\renewcommand{\headrule}{\hbox to\headwidth{\color{Couleur8}\leaders\hrule height \headrulewidth\hfill}}

% Environnements personnalisés
\newtcolorbox{seance}[1]{
    colback=Couleur6!10,
    colframe=Couleur1!65,
    fonttitle=\bfseries\large,
    title=Correction Séance \theseancenum #1,
    left=5mm,
    right=5mm,
    top=3mm,
    bottom=3mm,
    arc=0mm,
    after=\stepcounter{seancenum}
}

\newtcolorbox{seanceD}[1]{
    colback=Couleur6!5,
    colframe=Couleur6!45,
    fonttitle=\bfseries\large,
    title=Correction Devoirs - Séance \theseancenumD #1,
    left=5mm,
    right=5mm,
    top=3mm,
    bottom=3mm,
    arc=0mm,
    after=\stepcounter{seancenumD}
}

\begin{document}

\setcounter{seancenum}{25}
\newpage
\begin{seance}{} %25

\begin{enumerate}
\item $x-\dfrac{x}{5}=\dfrac{x\times 5}{5}-\dfrac{x}{5}=\dfrac{5x}{5}-\dfrac{x}{5}=\dfrac{5x-x}{5}=\dfrac{4x}{5}=$ ${\color[HTML]{f15929}\boldsymbol{\dfrac{4}{5}x}}$

\item On utilise l'égalité remarquable ${\color{red} a}^2-{\color{blue} b}^2=({\color{red} a}-{\color{blue} b})({\color{red} a}+{\color{blue} b})$ avec $a={\color{red} x}$  et $b={\color{blue} \sqrt{7}}$.\\$\begin{aligned}
         x^2-7&=\underbrace{{\color{red} x}^2-({\color{blue} \sqrt{7}})^2}_{a^2-b^2}\\
         &=\underbrace{({\color{red} x}-{\color{blue} \sqrt{7}})({\color{red} x}+{\color{blue} \sqrt{7}})}_{(a-b)(a+b)}
         \end{aligned}$ 

\item  Diminuer de $22\,\%$ revient à multiplier par $\dfrac{78}{100}$.\\
Si $V_I$ est le prix initial, on a : $V_I \times \dfrac{78}{100}=132$.\\
Ainsi, le prix initial est donné par : ${\color[HTML]{f15929}\boldsymbol{132 \times \dfrac{100}{78}}}$.\\La bonne réponse est la réponse {\bfseries \color[HTML]{f15929}A}.

\item $3\,000$ $\text{ m}$ $= 3$ $\text{km}$\\
L'aire du carré est : $3 \text{ km}\times 3 \text{ km} = {\color[HTML]{f15929}\boldsymbol{9\text{ km}^2}}$ .

\item  On a $4=2^{2}$ donc :\\
    $
    \text{Le double de  } 4^{50}  = 2 \times 4^{50} 
     = 2\times  \left(2^{2}\right)^{50} 
    =2\times  2^{100} 
    ={\color[HTML]{f15929}\boldsymbol{2^{101}}} 
$

\item $\dfrac{x-4}{x+12} \leqslant 0$\\
\medskip

$\bullet$ On commence par chercher les éventuelles valeurs interdites :\\$x +12 = 0$ \\$x +12 {\color[HTML]{f15929}\boldsymbol{-12}} = {\color[HTML]{f15929}\boldsymbol{-12}}$\\$x = -12$\\Le quotient est défini sur $\mathbb{R} \smallsetminus \{-12\}$.\\$\bullet$ On résout l'inéquation sur $\mathbb{R} \smallsetminus \{-12\}$.\\$x -4 > 0$ si et seulement si $x > 4$.\\$x+12 > 0$ si et seulement si $x > -12$.\\On peut donc en déduire le tableau de signes suivant :\\\begin{tikzpicture}[baseline, scale=0.75]
    \tkzTabInit[lgt=8,deltacl=0.8,espcl=3.5]{ $x$ / 2, $x-4$ / 2, $x+12$ / 2, $\dfrac{x-4}{x+12}$ / 4}{ $-\infty$, $-12$, $4$, $+\infty$}
	\tkzTabLine{  , -, t, -, z, +}
	\tkzTabLine{  , -, z, +, t, +}
	\tkzTabLine{  , +, d, -, z, +}
	\end{tikzpicture}
\\ L'ensemble de solutions de l'inéquation est $S = \left] -12 ; 4 \right] $.


\end{enumerate}
\end{seance}

\end{document}